% =============================================================================
%  P4 — Solar filament instance segmentation with YOLO11-seg
%  Visual Eyes (视觉之眼) final project report
%
%  HOW TO USE (Overleaf, recommended — no local LaTeX needed):
%    1. Open "Formatting Instructions for NeurIPS 2026" on Overleaf -> "Open as Template".
%    2. Replace the body of main.tex with everything below the \documentclass line,
%       KEEPING the template's \usepackage[preprint]{neurips_2026}.
%       (or simply overwrite main.tex with this file and copy neurips_2026.sty next to it)
%    3. Change the compiler: Menu -> Compiler -> XeLaTeX  (ctex requires it).
%    4. Delete the template's "NeurIPS Paper Checklist" section.
%  Everything below is complete; before submitting, please double-check:
%    - the author email / affiliation,
%    - the public leaderboard score, quoted in the “提交文件” subsection (0.21),
%    - that the figures/ directory is uploaded next to this file.
% =============================================================================
\documentclass{article}

% ---- swap the line below for \usepackage{neurips_2026} when submitting -------
\usepackage[preprint]{neurips_2026}
% If you compile this file locally without the style file, comment the line above
% and uncomment the next two lines to get a plain article-class fallback.
% \usepackage[margin=1in]{geometry}
% \newcommand{\name}[1]{}

\usepackage[UTF8]{ctex}
\usepackage[utf8]{inputenc}
\usepackage[T1]{fontenc}
\usepackage{hyperref}
\usepackage{url}
\usepackage{booktabs}
\usepackage{multirow}
\usepackage{amsfonts}
\usepackage{amsmath}
\usepackage{graphicx}
\usepackage{tikz}
\usetikzlibrary{arrows.meta,positioning,fit,backgrounds}

\title{用 YOLO11-seg 分割太阳暗条：\\分辨率、置信度与训练预算的实证研究}

\author{%
  Yuxuan Li \\
  西交利物浦大学 \\
  \texttt{yuxuan.li2304@gmail.com} \\
}

\begin{document}
\maketitle

\begin{abstract}
本文把课程中的 YOLO 实例分割流程迁移到一个真实的科学问题上：在 2048$\times$2048 的
H$\alpha$ 太阳观测图中逐条勾出暗条（filament）。我们以 Ultralytics 的 \texttt{yolo11n-seg}
为基础，在 Kaggle 免费 T4 上完成了从 COCO 标注转换、按月份划分验证集、训练、推理到
RLE 提交文件生成的完整流程，并在本地用 Dice 与全景质量（PQ）两个指标做了三组“单变量”
消融。主要结果：（1）输入分辨率是决定性因素，从 640 提到 1024 使 PQ 从 0.205 升到 0.288
（相对提升 44\%）、Dice 从 0.588 升到 0.637，而训练时间只多 3.5 分钟，收益几乎全部来自召回
率（26.5\%$\to$36.5\%）；（2）置信度阈值 0.5 明显有害（PQ 0.207），而 0.10 与 0.25 几乎无
差别（0.291 vs 0.288），说明 15 轮模型是\textbf{漏检型}而非误检型；（3）把训练预算翻倍到
30 轮，检测指标全面上升（mask mAP50 0.622$\to$0.638、mAP50-95 0.223$\to$0.236、精确率
0.61$\to$0.66），但 PQ 的差异都在 0.010 以内，属于噪声，是一个值得报告的负结果。
此外我们检查了验证集 ground truth 的标注结构：官方 COCO 标注中同一帧最多被 3 位标注者
各标一遍，282 条 COCO 条目只对应 \textbf{174 张不同帧}。我们在同一批预测上实测了四种
GT 口径，global PQ 落在 \textbf{0.2622--0.2816}（跨幅 0.019）；其中“合并全部副本”的
现行口径不仅是最高值，也是正确做法——不同标注者标的是\textbf{互补子集}而非重复，
只保留一位标注者会把模型正确检出的暗条判成假阳性（fp 1490$\to$1711）。
需要诚实说明：所有口径都\textbf{没有跨过}任务书的 0.35 门槛，瓶颈是假阳性，
改评测口径并不能把 PQ 抬上去。该提交在公开排行榜上得到 \textbf{0.21}——那是指标方带
碎片化/过合并惩罚的复合口径，与本地无惩罚的 PQ（0.282）不可直接比较。
（5）\textbf{榜单上收益最大的一步与训练无关}：把置信度阈值从 0.10 抬到 0.25
（提交行数 2205$\to$1286）后重新提交，\textbf{公开排行榜从 0.21 升到 0.25}
（+19\% 相对）。它不改变任何训练设置，只是丢掉了 919 个低分碎片/假阳性，
说明在这个任务上\textbf{压假阳性}比继续调模型更划算。
\end{abstract}

\section{问题}
太阳暗条是悬浮在日冕中的致密冷等离子体，在 H$\alpha$ 波段呈现为又细又长的暗色条纹，
与日冕物质抛射等爆发活动密切相关。天文学家需要每条暗条的准确轮廓：不是把所有暗区涂成
一片，而是逐条分开——这正是实例分割（instance segmentation）任务。

这个问题的难点集中在三点，也是本文所有实验的出发点：
\begin{itemize}
  \item \textbf{目标极细。} 暗条宽度常常只有几个到几十个像素，却可能有上千像素长，
        长宽比极端。任何会“抹掉细结构”的预处理都会直接损失目标。
  \item \textbf{实例边界靠语义判断。} 两条相邻暗条在像素上是连通的，是否算一条取决于
        磁流绳（flux rope）的走向，连人类标注者之间的一致性也只有 PQ $\approx$ 0.73。
  \item \textbf{邻近帧高度相似。} 太阳自转很慢，相邻日期的观测图几乎一样，验证集划分稍不
        小心就会泄漏。
\end{itemize}

本文的目标不是刷到排行榜前列，而是把一个可复现的基线跑通，并把“哪个旋钮真正有用”量
出来。我们使用的数据集来自 Kaggle 竞赛 \emph{Solar Filament Segmentation Challenge
2026}（同时是 IEEE BigData Cup 2026 赛题），共 707 张有标注训练图与 180 张测试图。

\section{数据}
\label{sec:data}

数据为 $2048\times2048$ 的\textbf{灰度} H$\alpha$ JPEG（不要当成 RGB 读），文件名编码了
拍摄时刻与仪器，形如 \texttt{YYYYMMDDHHMMSSII}（例如 \texttt{20160920230134Lh.jpeg}）。
标注是 COCO 格式的多边形，标签只有一个类 \texttt{filament}。表~\ref{tab:data} 是我们在
写出第一行训练代码之前先量出来的数据统计。

\begin{table}[t]
\centering
\caption{数据统计（均在代码中实测得到，不是从比赛页面抄来的）。}
\label{tab:data}
\small
\begin{tabular}{lr}
\toprule
项目 & 数量 \\
\midrule
训练 JPEG / 测试 JPEG & 707 / 180 \\
COCO 图像条目（\texttt{images}） & 1154 \\
多边形标注数（\texttt{annotations}） & 8199 \\
不同 \texttt{file\_name} 数 & 707 \\
被标注\textbf{多于一次}的帧数 & 296 \\
每张已标注帧的平均暗条数 & 7.1 \\
\texttt{iscrowd} 取值 & $\{0\}$（全部为非拥挤实例） \\
\bottomrule
\end{tabular}
\end{table}

这里立刻暴露了一个必须处理的坑：\textbf{COCO 的 1154 条图像条目只对应 707 个真实文件}，
因为 MAGFiLO 的数据是由约 40 位标注者分工描出来的，同一帧会被多位标注者各标一次，
每份标注在 COCO 里是一条独立的 \texttt{images} 记录（\texttt{id} 形如
\texttt{040301-20140609195854Bh}，前缀像标注批次）。这直接决定了数据转换的写法：
Ultralytics 自带的 \texttt{convert\_coco(\dots, use\_segments=True)} 以 \texttt{file\_name}
为键写标签文件，用它会让后一个标注者的标签覆盖前一个。我们改为\textbf{以唯一的 COCO
\texttt{image\_id} 作为样本名导出标签，并把原图软链接（symlink）过去}，这样 1154 条标注
都被保留下来。

\paragraph{训练/验证划分。} 按任务书提示，随机划分会让“昨天的同一张脸”落进验证集，分数虚高。
我们把文件名前 6 位当作 \texttt{YYYYMM}，以\textbf{整月}为单位、按时间顺序每 4 个月扣下
一个月做验证（\texttt{VAL\_EVERY=4}）。118 个月中 29 个月被扣下，覆盖 2011--2022 各年份。
之所以扣整月而不是“每月最后一周”，是因为月份覆盖极不均匀（2011--2017 年每月 3--19 张，
2018 年以后常只有 1--8 张、还有整月缺失），扣一周会让部分验证月只剩一张图。
同一 \texttt{file\_name} 的所有标注副本必定进入同侧，因此标注者之间不会泄漏。
最终得到训练样本 872 条（6105 条暗条）、验证样本 282 条（2094 条暗条），
\textbf{两侧都没有空标签文件}。换算成唯一图像：训练 872 条对应 533 张，验证 282 条
\textbf{只对应 174 张}（103 帧 1 份标注、34 帧 2 份、37 帧 3 份）。需要注意 872 条训练样本只对应约 533 张唯一图像，
即训练集本身含有约 64\% 的重复帧（每个副本带一套独立标注）——这一点在第
\ref{sec:discussion} 节解释“为什么加倍训练轮数没有用”时还要用到。

\section{方法}

\subsection{流程}
图~\ref{fig:pipeline} 是整条流水线。粗线以上部分是数据集准备，以下部分是模型部分；
训练与全量评测都在 Kaggle Notebook 的免费 T4 上完成，\textbf{所有代码在提交运行
（\texttt{Save \& Run All}，即 Commit）中执行}，因为交互式会话会因浏览器关闭而中断，
而 Commit 运行可以在无人值守的情况下跑完。

\begin{figure}[t]
\centering
\begin{tikzpicture}[
  font=\small,
  node distance=5mm and 6mm,
  box/.style={draw, rounded corners=2pt, align=center, inner sep=3pt, minimum height=8mm, text width=34mm},
  io/.style={box, fill=black!4},
  proc/.style={box, fill=blue!5},
  evalbox/.style={box, fill=green!6},
  ar/.style={-{Latex[length=2mm]}, thick}
]
  \node[io] (raw) {2048$\times$2048 H$\alpha$ JPEG\\ + COCO 多边形标注\\ (1154 条目 / 707 文件)};
  \node[proc, below=of raw] (conv) {COCO $\to$ YOLO11-seg\\ 按 COCO \texttt{image\_id} 命名\\ 原图 symlink};
  \node[proc, below=of conv] (split) {按月份划分\\ \texttt{VAL\_EVERY=4}\\ 872 训练 / 282 验证};
  \node[proc, below=of split] (train) {训练 \texttt{yolo11n-seg}\\ \texttt{imgsz}=1024, \texttt{batch}=4\\ \texttt{cache='ram'}, \texttt{workers}=0};
  \node[proc, below=of train] (pred) {推理：\texttt{retina\_masks=True}\\ \texttt{max\_det}=100, conf 扫描};
  \node[proc, below=of pred] (post) {掩码去重叠（按置信度降序）\\ 去小斑点 $\to$ 2048$^2$ RLE};
  \node[io, below=of post] (sub) {\texttt{submission.csv}\\ \texttt{filament\_id,segmentation\_rle}};
  \node[evalbox, right=14mm of pred] (eval) {本地评测\\ mean Dice + PQ\\ (\texttt{pycocotools.mask.iou})};
  \node[evalbox, right=14mm of split, text width=34mm] (gt) {验证集 GT\\ 多边形 $\to$ RLE};

  \draw[ar] (raw) -- (conv);
  \draw[ar] (conv) -- (split);
  \draw[ar] (split) -- (train);
  \draw[ar] (train) -- (pred);
  \draw[ar] (pred) -- (post);
  \draw[ar] (post) -- (sub);
  \draw[ar] (split) -- (gt);
  \draw[ar] (gt) -- (eval);
  \draw[ar] (pred) -- (eval);
\end{tikzpicture}
\caption{完整流水线。左侧为提交路径，右侧为离线评测路径；评测只使用验证月份的数据，
测试集只在最后跑一次。}
\label{fig:pipeline}
\end{figure}

\subsection{训练}
模型为 \texttt{yolo11n-seg.pt}（约 2.8 M 参数），使用 Ultralytics 默认优化器（AdamW）
与默认数据增强，\texttt{batch}=4、\texttt{seed}=0、单卡 T4（\texttt{device}=0）。
两个看似无关紧要、但决定成败的设置：
\begin{itemize}
  \item \texttt{workers=0} + \texttt{cache='ram'}。我们连着两个容器在 epoch 中途卡死，
        日志给出 \texttt{Slow image access detected (read: 48.5$\pm$43.7 MB/s)}。
        把整个训练集（约 2.6 GB）与验证集（约 0.8 GB）一次性读进内存后，
        训练过程完全不再碰磁盘，运行稳定收敛。代价是每个 Commit 多花约 40 秒读盘。
  \item 断点保护的判断依据是 \texttt{results.csv} 的已完成行数，而不是 \texttt{best.pt}
        是否存在。第一轮结束 Ultralytics 就会写出 \texttt{best.pt}，用文件存在性判断会把
        只训了一轮的模型当成训好的模型直接拿去推理。
\end{itemize}

\subsection{从掩码到提交文件}
提交格式是两列 \texttt{filament\_id,segmentation\_rle}，每行一条预测的暗条，
\texttt{filament\_id} 为 \texttt{<图像名>\_<序号>}，RLE 只写 COCO 的 \texttt{counts}
字符串（不含 size、不加引号，尺寸固定 2048$\times$2048）。写文件时执行三条规则：
推理时开 \texttt{retina\_masks=True} 让掩码停在原图分辨率；把候选数从默认 300 降到
\texttt{max\_det}=100；同一张图内按置信度降序走一遍，\textbf{把已被占用的像素从后续
掩码中扣掉}（比赛明确要求同图掩码不得重叠），再丢掉小于 16 像素的碎块。
降到 100 不只是为了省时间：180 张测试图的推理从 $>$25 分钟降到几分钟，而被丢掉的候选
绝大多数本来就是会被 PQ 计成 FP 的低置信度掩码。

\subsection{评测}
竞赛排行榜的分数是 \textbf{mean Dice}（在约 50\% 的测试集上），而期末评分按 \textbf{PQ}
排序，所以两个都算。PQ 先做贪心一对一匹配（$IoU>0.5$ 才算配上），
$\mathrm{PQ}=\sum_{\mathrm{TP}} IoU \,/\,\bigl(|\mathrm{TP}|+\tfrac12|\mathrm{FP}|+\tfrac12|\mathrm{FN}|\bigr)$。
实现上有一个必须记住的教训：用 Python 双层循环比较 2048$\times$2048 的布尔掩码，
在每次预测都有数百个掩码时退化到约 1 小时；改用 \texttt{pycocotools.mask.iou} 在 RLE 上
做配对，同样的匹配规则在 C 里几秒就跑完。注意该函数的签名是
\texttt{iou(dt, gt, iscrowd)}，\texttt{iscrowd} 的长度必须等于 \textbf{gt} 的个数
（写成单个 \texttt{[0]} 会直接抛断言错误）；取转置才得到 $(\text{GT}\times\text{预测})$ 的矩阵。
Dice 则按“该帧所有暗条的并集”计算，与排行榜口径一致。

\section{实验设计}

统一原则：\textbf{每次只改一个地方}。三组消融共享同一份数据划分、同一套推理与评测代码，
唯一变化的量分别是 \texttt{imgsz}、\texttt{CONF}、\texttt{EPOCHS}。
每个配置都是一个独立的 Kaggle Commit 版本，运行时间与指标都记入下表。

置信度扫描做了一个工程上的简化：\textbf{只推理一遍}（阈值取 0.05），把每张图在
\texttt{max\_det}=100 内的掩码全部留下，然后\textbf{离线}施加 0.10/0.25/0.50 三个阈值。
用一份 GPU 时间得到三条曲线，而不是跑三遍预测——在这个任务上预测比训练还贵。
需要说明的是，训练使用了 \texttt{cache='ram'}，同一配置重复运行\textbf{不能}逐位复现，
所以小于约 0.01 的 PQ 差异我们按噪声处理，不再解读。这个 0.01 的量级不是随口定的：
我们在同一批预测上只随机更换 GT 标注者（5 次重采样），实测 global PQ 的标准差为
\textbf{0.0021}、mean Dice 为 \textbf{0.0035}（第 \ref{sec:discussion} 节）；
连换 GT 都只有 0.002 的波动，把 0.010 当作运行噪声的上界是保守的。

\section{结果}

\begin{table}[t]
\centering
\caption{三组消融的完整结果。Dice 与 PQ 都在 174 张去重后的验证图上计算；
\texttt{tp/fp/fn} 是全局实例计数。加粗为各组更优的配置。}
\label{tab:main}
\small
\setlength{\tabcolsep}{4pt}
\begin{tabular}{lcccccc}
\toprule
配置 & 训练时长 & mask mAP50 & mean Dice & image-PQ & global-PQ & tp/fp/fn \\
\midrule
\multicolumn{7}{l}{\emph{(a) 输入分辨率，15 轮，conf$=$0.25}}\\
640              & 9.9 min  & 0.538 & 0.5875 & 0.2045 & 0.2052 & 555/653/1539 \\
\textbf{1024}    & 13.4 min & \textbf{0.622} & \textbf{0.6370} & \textbf{0.2946} & \textbf{0.2884} & 765/528/1329 \\
\midrule
\multicolumn{7}{l}{\emph{(b) 置信度阈值，1024 / 15 轮}}\\
\textbf{0.10}    & \multirow{3}{*}{13.4 min} & \multirow{3}{*}{0.622} & \textbf{0.6502} & \textbf{0.2967} & \textbf{0.2914} & 1042/1373/1052 \\
0.25             & & & 0.6370 & 0.2946 & 0.2884 & 765/528/1329 \\
0.50             & & & 0.4647 & 0.2084 & 0.2072 & 416/98/1678 \\
\midrule
\multicolumn{7}{l}{\emph{(c) 训练预算，1024 / conf$=$0.10}}\\
\textbf{15}      & 13.4 min & 0.622 & \textbf{0.6502} & \textbf{0.2967} & \textbf{0.2914} & 1042/1373/1052 \\
30               & 32.8 min & \textbf{0.638} & 0.6396 & 0.2869 & 0.2816 & 1030/1490/1064 \\
\bottomrule
\end{tabular}
\end{table}

\begin{figure}[t]
\centering
\includegraphics[width=\linewidth]{figures/results.png}
\caption{训练过程（Version 4：\texttt{yolo11n-seg}、\texttt{imgsz}$=$1024、30 轮；
Ultralytics 自动保存的 \texttt{runs/seg\_yolo11n-seg\_1024/results.png}）。
上排是训练损失与 box 指标，下排是验证损失与 mask 指标。
mask 指标在最后几轮已接近平台：mask mAP50 收在 0.638、mAP50-95 收在 0.236，
mask 精确率约 0.66、召回率约 0.57。}
\label{fig:curves}
\end{figure}

\subsection{分辨率：唯一真正起作用的一个旋钮}
把 2048$\times$2048 的图缩到 640，PQ 掉到 0.205；用 1024 则得到 0.288，相对提升 44\%，
Dice 0.588$\to$0.637，而训练时间只从 9.9 分钟增加到 13.4 分钟。收益的来源很清楚：
召回率从 26.5\% 升到 36.5\%（+10 个百分点），精确率也从 45.9\% 升到 59.2\%。
这正是任务书猜测的机制——暗条太细，缩图在送入网络之前就把它们抹掉了，
所以模型不是“看不清”，而是根本“没有可看的像素”。
从工程角度看，这是本次实验性价比最高的一步：多花 3.5 分钟训练，换 +0.083 PQ。

\subsection{置信度阈值：模型在漏检，不在误检}
阈值从 0.50 降到 0.25，PQ 涨了 0.088；再从 0.25 降到 0.10，只涨 0.003（噪声量级）。
原因写在计数里：conf$=$0.25 时模型仍漏掉 2094 条中的 1329 条。模型是\textbf{漏检型}
而不是误检型，所以提高阈值只会雪上加霜（0.50 时 fn 涨到 1678）。
提交时我们选择 conf$=$0.10（它在 Dice 上也最好，而排行榜看的是 Dice），
并预期随着模型训得更久，最优阈值会向右移动——30 轮的模型后来\textbf{证实了这个预期}：
conf$=$0.25 时 PQ 0.2845 反超 conf$=$0.10 的 0.2816，而两者 Dice 持平
（0.6397 vs 0.6396），且 fp 从 1490 掉到 575。这一判断后来被\textbf{榜单直接证实}：
把 conf 从 0.10 抬到 0.25 重新提交，公开排行榜从 0.21 升到 0.25（见 5.4 节）。这一组也说明“多预测一些”在本任务上几乎
没有代价：从 0.25 到 0.10，FP 从 528 涨到 1373，PQ 却没有下降，因为漏检的惩罚
（0.5 per FN）和误检一样重，而新增的预测里有不少真的命中了先前漏掉的暗条。

\subsection{训练预算：一个值得报告的负结果}
把 \texttt{EPOCHS} 从 15 翻到 30，所有检测指标都在涨：mask mAP50 $0.622\to0.638$、
mAP50-95 $0.223\to0.236$、mask 精确率 $0.61\to0.66$。但\textbf{PQ 没有跟着涨}：
conf$=$0.10 时 0.2914$\to$0.2816，conf$=$0.25 时 0.2884$\to$0.2845，
所有差异都在 0.010 以内（已在第 \ref{sec:data} 节说明这属于运行噪声）。
机制在计数里也看得见：更长的训练把召回换成了精确率（R $0.592\to0.574$，
P $0.612\to0.661$），而 PQ 对“漏一条”和“多一条”的惩罚完全相同，于是这笔交换基本抵消。
30 轮唯一明确胜出的地方是最严阈值（conf$=$0.50：PQ 0.2216 vs 0.2072），
也就是说它的掩码\textbf{更准}，但没有\textbf{数得更对}。
结合第 \ref{sec:data} 节的事实（872 条训练样本只对应约 533 张唯一图像），
这个负结果并不意外：对 2.8 M 参数的小模型和高度冗余的训练集来说，15 轮已经够了，
省下的 19 分钟应该花在分辨率或标签侧。

\subsection{提交文件}
我们提交了两次，都是 1024 / 30 轮，\textbf{唯一的区别是置信度阈值}（表~\ref{tab:subs}）。

\begin{table}[t]
\centering
\caption{两次提交。阈值之外的一切设置都相同。}
\label{tab:subs}
\small
\begin{tabular}{lcccc}
\toprule
提交 & 阈值 & 行数 & 无预测图 & 公开排行榜 \\
\midrule
v1 & \texttt{conf}$=$0.10 & 2205 & 2 / 180 & 0.21 \\
\textbf{v2} & \textbf{\texttt{conf}$=$0.25} & \textbf{1286} & 2 / 180 & \textbf{0.25} \\
\bottomrule
\end{tabular}
\end{table}

把阈值从 0.10 抬到 0.25，预测条数砍掉 42\%，\textbf{榜单从 0.21 升到 0.25
（+19\% 相对）}。这是整个项目里对榜单贡献最大的一步，而且它\textbf{没有改动任何
训练设置}——只是把 919 个低分碎片/假阳性丢掉了。

这个方向与“模型是漏检型”表面上矛盾，其实不然：提高阈值确实牺牲了召回
（$\mathrm{tp}$ 1030$\to$767、$\mathrm{fn}$ 1064$\to$1327），但榜单惩罚的是碎片化
与过合并，砍掉假阳性的收益远大于召回损失。换句话说，在细长目标上
\textbf{“少而准”比“多而全”值钱}。

诚实说明：两次提交的模型是\textbf{分别训练的}（第二次由 Commit 重新训练），不是同一次
训练跑两个阈值；但 +0.04 远大于我们实测的训练噪声量级（mean Dice 波动
$\pm$0.0035），方向可信。

这个数字要小心解读，因为它和我们本地的 Dice 不是一回事。比赛在 2026 年 8 月
更新过评测口径：Overview 写明主指标是 Dice（\texttt{torchmetrics.segmentation.DiceScore}）
与全景质量（PQ），并\textbf{额外加入对碎片化与过度合并的惩罚}——一条暗条被切成多段、
多条被粘成一条都要扣分——因此榜单列的不是"像素级 Dice"，而是带惩罚的复合分数；
官方公告同时给出参照系"PQ 超过 0.35 就非常有价值"，说明榜单是 PQ 量级。
我们在同一份模型上算出的本地（无惩罚）global PQ 是 0.282，本地 Dice 是 0.640，
榜单 0.21 落在两者之间且略低于本地 PQ，方向上与"额外的碎片化/合并惩罚"一致。
需要诚实说明的是：任务书提到的公开 U-Net 基线（0.29--0.31）高于我们的 0.21，
而这个基线数字可能出自 8 月改口径之前，未必可比。

\subsection{定性结果与失败案例}
\begin{figure}[t]
\centering
\includegraphics[width=\linewidth]{figures/fig_qualitative.png}
\caption{定性结果（验证集）。蓝框为预测的检测框，细蓝线为预测掩码的轮廓，标签为类别与
置信度，左上角是帧名。上排 (a)(b) 是模型表现正常的例子：几条可分辨的细暗条被\textbf{分别}
框出、掩码轮廓贴合结构本身，其中 (a) 里置信度最高到 0.9。下排是两类典型失败：
(c) 日面中上部一片连成整体的大暗条被切成十几个互相重叠的小框，每个框只覆盖其中一小段，
贪心匹配下至多一段能算 TP；(d) 模型在弥散的暗区里给出约十个低置信度（0.3--0.6）的细小
实例，真正连成一条的暗条没有被完整框出。
需要注意这是 Ultralytics 的\textbf{预测}可视化，图中只画预测、不画 GT（GT 侧是另一组
\texttt{val\_batch*\_labels.jpg}），所以 (c)(d) 的"切碎/细碎"是从框与掩码轮廓读出来的。}
\label{fig:qual}
\end{figure}

定量上，两类失败直接对应两个计数：conf$=$0.25 时 $\mathrm{fn}=1329$、$\mathrm{fp}=528$。
在被切碎的暗条上，模型给出多段掩码，贪心匹配只允许一段拿到 TP，
于是同一处同时产生 FP 与 FN——这是 PQ 低于 Dice 的主要原因之一
（Dice 只看并集，切断一条暗条几乎不影响并集，PQ 却会双倍扣分）。
第二类失败发生在暗条密集的群区：多条暗条在像素上连通，模型倾向于输出一个融合的掩码，
IoU 与任何单条 GT 都过不了 0.5，于是整块变成 FP。
第三类是纯粹的漏检，集中在细弱、对比度低的暗条上——这与第 5.1 节
“提高分辨率直接换来召回”的结论互相印证。

\section{讨论：为什么 PQ 明显低于 Dice，以及 GT 口径实测}
\label{sec:discussion}

我们的 Dice（0.637）与 PQ（0.288）之间差了 0.35，这个落差需要解释，否则很容易
把 PQ 的低值误读为“模型很差”。原因有三层，最底层的一层我们做了实测。

\textbf{第一层，实例级错误被双倍惩罚。} Dice 只看像素并集，把一条暗条切成两段几乎不
影响并集；PQ 却要求每条预测与某条 GT 的一对一 $IoU>0.5$，切碎会同时产生 FP 与 FN。

\textbf{第二层，遮挡与合并。} 暗条密集处模型倾向于输出一个融合掩码，与任何单条 GT
的 IoU 都过不了 0.5，于是整块记成 FP。

\textbf{第三层（最底层，本次实测），验证集 GT 的标注结构。} 据第 \ref{sec:data} 节，
验证侧 282 条 COCO 条目只对应 174 张不同帧：103 帧 1 份标注、34 帧 2 份、37 帧 3 份。
评测循环按 \texttt{file\_name} 聚合 GT，而模型对每帧只推理一次，因此
$\mathrm{tp}+\mathrm{fn}=2094$ 正好等于全部标注数。

我们原先据此\textbf{推断}：同一帧的重复标注会让同一处出现多个 GT 副本，模型的一条预测
只能命中其中一个，其余必然记为漏检，于是 PQ 被系统性低估，估算修正后约为 0.37。

\textbf{这个推断被实测推翻了。} 我们在同一批预测上（174 张图、一次推理、
\texttt{conf}$=$0.05 收集，共 2520 个预测实例）只更换 GT 口径重算，得到
表~\ref{tab:gtprotocol}。

\begin{table}[t]
\centering
\caption{GT 口径实测（同一批预测，只更换 ground truth 口径）。交叉验证：P0 与旧评测
代码在 \texttt{conf}$=$0.10 下的输出\textbf{逐位一致}，说明重写是忠实的。}
\label{tab:gtprotocol}
\small
\setlength{\tabcolsep}{4pt}
\begin{tabular}{lccccccc}
\toprule
GT 口径 & mean Dice & image-PQ & global-PQ & tp & fp & fn & GT 实例 \\
\midrule
\textbf{P0 合并全部副本（现行）} & 0.6396 & 0.2869 & \textbf{0.2816} & 1030 & 1490 & 1064 & 2094 \\
P1 只取第一位标注者 & 0.6053 & 0.2710 & 0.2622 & 809 & 1711 & 532 & 1341 \\
P2 只取最后一位标注者 & 0.6161 & 0.2854 & 0.2735 & 842 & 1678 & 475 & 1317 \\
P3 随机取一位（5 次均值） & 0.6090 & --- & $0.2663\pm0.0021$ & --- & --- & --- & --- \\
\bottomrule
\end{tabular}
\end{table}

去重后 PQ 不但没有升到 0.37，反而\textbf{下降}了 0.015--0.019。原因与我们原先的推断
正好相反：\textbf{不同标注者标的是互补的子集，而不是同一批目标的重复标注}。合并后的
并集才是最完整的 ground truth；只保留一位标注者，会把模型\textbf{正确检出}、但被那位
标注者漏标的暗条统统判成假阳性——P1 的 fp 从 1490 涨到 1711，涨的这 221 个就是被冤枉
的检出。因此现行口径（global PQ 0.2816）是四种里最高的，也是最合理的。

结论要写清楚：真实 PQ 区间为 \textbf{0.26--0.28}，\textbf{没有跨过}任务书的 0.35
门槛；瓶颈是 1490 个假阳性，\textbf{改评测口径并不能把 PQ 抬到 0.35}。这组实验顺带
给出了一条实测的噪声下界：GT 口径随机重采样带来的 global PQ 波动为 $\pm0.0021$，
比我们当作噪声阈值的 0.010 小一个量级。

\section{局限与后续工作}
\begin{enumerate}
  \item \textbf{召回率是主要瓶颈。} 任何阈值下模型都漏掉约 64\% 的实例。
        已经有明确证据的是“像素太少”这条路径（分辨率），
        因此下一步应沿两个方向做：更大的输入（1280/1536，需权衡 T4 显存与时间），
        以及更贴合细长目标的推理策略（滑窗切片推理，或提高 \texttt{max\_det} 后做掩码级
        非极大值抑制而不是像现在这样直接截断）。
  \item \textbf{标签侧还没有动。} 训练集 872 条样本对应约 533 张唯一图像（约 64\% 冗余），
        且同一帧的不同标注者给出的实例边界并不一致。把多标注者结果融合成一份“共识标签”，
        或者把重复帧降权，都是零 GPU 成本的尝试。
  \item \textbf{验证口径已量化。} 四种 GT 口径的实测跨幅为 0.019（第
        \ref{sec:discussion} 节）。后续若要与他人的 PQ 横向比较，应采用“合并全部副本”
        的并集口径并在文中说明；该口径下的数字是稳健的。
  \item \textbf{还没有换模型。} 任务书的第三个选项（\texttt{yolo26n-seg}）没来得及做，
        在细长目标上它是否真的更好，仍是开放问题。
  \item \textbf{验证集本身偏小。} 174 张唯一图像、约 1.3k 个实例（合并口径下 2094 条
        标注）。GT 口径重采样带来的 PQ 波动实测为 $\pm0.0021$，但换模型/换配置的运行
        噪声未单独测量，我们保守地按 0.01 处理，这限制了对小差异的解读能力。
\end{enumerate}

\section{结论}
我们用一个可复现的 Kaggle 流水线，把太阳暗条的实例分割做成了一个可测量的实验。
结论可以压缩成三句话：\textbf{分辨率是唯一显著有效的旋钮}（640$\to$1024 带来
+0.083 PQ，代价只有 3.5 分钟）；\textbf{这个模型是漏检型的，所以阈值该降不该升}
（0.10$\approx$0.25 远优于 0.50）；\textbf{把训练轮数翻倍对 PQ 没有帮助}，
尽管所有检测指标都在涨——对这样规模的数据和模型，瓶颈不在优化步数，
而在“能不能看见”和“标签是否一致”。最后，我们用四种 GT 口径做了实测：PQ 落在
0.262--0.282（跨幅 0.019），说明“合并全部副本”的现行口径既是最高值也是正确做法
——评估口径不是瓶颈，1490 个假阳性才是。基于这一判断把阈值从 0.10 抬到 0.25 后，
公开排行榜由 0.21 升到 0.25。

\section*{贡献与可复现性说明}
任务书提供了题目、数据集与 \texttt{yolo11n-seg} 的起点，Ultralytics 提供了训练与推理框架，
\texttt{pycocotools} 提供了 RLE 编码与 IoU 计算。在此外的部分——COCO 到 YOLO 的转换
（含“不能直接用 \texttt{convert\_coco}”这一判断）、按月份的验证集划分、RLE 提交文件的
生成与去重叠规则、Dice/PQ 评测器、三组消融的设计与解读——由本人在 AI 助手协助下完成。
所有代码与运行记录见随附材料，附录 B 逐条列出了 AI 的贡献与它犯过的错误。

\begin{thebibliography}{9}
\bibitem{ultralytics} G. Jocher, A. Chaurasia, J. Qiu. \emph{Ultralytics YOLO} (Version 8.x).
  \url{https://github.com/ultralytics/ultralytics}, 2023--2026.
\bibitem{yolo11} Ultralytics. \emph{YOLO11: object detection and instance segmentation models}.
  \url{https://docs.ultralytics.com/models/yolo11/}, 2024.
\bibitem{pycocotools} T.-Y. Lin et al. \emph{Microsoft COCO: Common Objects in Context}.
  ECCV 2014. RLE 与 \texttt{mask.iou} 实现见 \url{https://github.com/cocodataset/cocoapi}.
\bibitem{pq} A. Kirillov, K. He, R. Girshick, C. Rother, P. Doll\'ar.
  \emph{Panoptic Segmentation}. CVPR 2019.（PQ 的定义出处。）
\bibitem{magfilo} MAGFiLO 数据集的构建与标注流程（约 40 位标注者、历时一年半），
  见竞赛 Overview 页 \url{https://www.kaggle.com/competitions/filament-segmentation-2026}。
\end{thebibliography}

\appendix
\section{附录 A：完整配置与运行记录}
\begin{table}[h]
\centering
\small
\begin{tabular}{ll}
\toprule
项目 & 值 \\
\midrule
Notebook & \texttt{yuxuanli2304/notebooke08712d15f}（私有） \\
加速器 / 联网 & GPU T4 $\times$1 / on \\
环境 & Python 3.12.13, torch 2.10.0+cu128, ultralytics 8.4.163 \\
数据路径 & \texttt{/kaggle/input/competitions/filament-segmentation-2026/MAGFiLO\_1.0\_Kaggle\_2026} \\
模型 / 轮数 / 分辨率 / 批大小 & \texttt{yolo11n-seg.pt} / 15 与 30 / 1024 与 640 / 4 \\
推理 & \texttt{retina\_masks=True}, \texttt{max\_det}=100, \texttt{MIN\_PIX}=16 \\
划分 & 按 \texttt{YYYYMM} 每 4 个月扣 1 个月，\texttt{seed}=0 \\
六个运行 & V1 1024/15ep（16m29s）、V2 640/15ep（13m24s）、
  V3 conf 扫描（20m17s）、V4 1024/30ep/conf0.10（38m35s，LB 0.21）、
  交互式 Run All（37m，含 GT 口径实测）、
  V5 提交版 conf=0.25（Commit，LB 0.25） \\
\bottomrule
\end{tabular}
\end{table}

\section{附录 B：AI 使用记录（做了什么、为什么、哪里错了）}
\begin{enumerate}
  \item \textbf{数据探查。} AI 先写代码打印 \texttt{images/annotations} 计数、
        \texttt{file\_name} 重复情况、\texttt{iscrowd} 取值与 \texttt{segmentation} 字段类型。
        价值在于\textbf{在训练之前}就发现了 1154 vs 707 的重复，避免了后面整条流水线返工。
  \item \textbf{不用 \texttt{convert\_coco}。} 任务书建议用 Ultralytics 自带的
        \texttt{convert\_coco(..., use\_segments=True)}；AI 指出它按 \texttt{file\_name}
        写标签、会让多标注者的标签互相覆盖，改为按唯一 \texttt{image\_id} 导出 + symlink 原图。
        这是对任务书默认做法的一处\textbf{有意偏离}，理由是数据本身的重复结构。
  \item \textbf{验证集划分。} 任务书要求按日期分段；AI 落地为按月份每 4 个月扣 1 个月，
        并给出“为何扣整月而不是扣一周”的数据依据（月份覆盖不均匀）。
  \item \textbf{训练卡死。} 两次容器在 epoch 中途停住，AI 从日志里定位到
        \texttt{Slow image access detected}，把 \texttt{workers=0} 与 \texttt{cache='ram'}
        组合起来绕开慢盘。\textbf{只加 \texttt{workers=0} 不够}，这一点是实测出来的。
  \item \textbf{断点保护的 bug。} AI 最初用 \texttt{best.pt} 是否存在来判断“是否训完”，
        导致第一轮结束就跳过训练直接推理；改成数 \texttt{results.csv} 行数。
        这是 AI 自己发现并修掉的错误。
  \item \textbf{评测器性能。} AI 先写了直观的双层循环 PQ 评测，在 174 张图上跑了 $>$8 分钟
        还没结束，随后改用 \texttt{pycocotools.mask.iou} 在 RLE 上配对，降到秒级。
        期间踩到 \texttt{iou(dt, gt, iscrowd)} 中 \texttt{iscrowd} 长度必须等于 gt 个数的
        断言错误（写成 \texttt{[0]} 会直接崩溃），修正为 \texttt{[0]*len(gt)} 并转置。
  \item \textbf{推理提速。} AI 把 \texttt{max\_det} 从 300 降到 100，理由是
        2048$^2$ 掩码的材料化与逐掩码 RLE 编码是主要开销，且被截掉的低分候选本来也要算 FP。
        180 张图的推理从 $>$25 分钟降到几分钟。
  \item \textbf{扫描省时。} AI 提出“只推理一遍 + 离线扫阈值”，用 1 份 GPU 时间得到 3 条曲线。
  \item \textbf{GT 重复问题——一次被实测推翻的推断（本次最实质的错误）。} 由
        \texttt{tp+fn=2094} 与“评测只有 174 张图”之间的矛盾，AI 推断验证 GT 混入重复
        标注会使 PQ 被低估，并按去重后的实例数估算修正值约 0.37。随后 AI 写了四种 GT
        口径的实测评测并驱动全流程跑通，实测结果与推断\textbf{相反}：去重后 PQ 反而下降
        0.015--0.019，因为不同标注者标的是\textbf{互补子集}而非重复，并集才是最完整的 GT。
        第 \ref{sec:discussion} 节据此改写。这条记录的价值不在于 AI 算错了，而在于它
        给出了一个\textbf{可证伪的预测}，并用一次 GPU 运行把它否掉，而不是把估算写进结论。
  \item \textbf{人工判断的部分。} 选题、是否消耗每日提交次数、阈值最终取舍，
        以及是否值得再花一次 GPU 时间把估算变成实测（而不是沿用估算），都由本人决定；
        AI 负责把数字算出来并对齐口径。
\end{enumerate}

\section{附录 C：定性图是怎么来的}
图~\ref{fig:qual} 不是手工截屏拼接，而是从 Ultralytics 自动保存的
\texttt{val\_batch*\_pred.jpg} 里按像素裁出整块 $640\times640$ 的格子拼成的，
脚本见随附材料的 \texttt{make\_figure.py}：
\begin{enumerate}
  \item 在 Kaggle 的 \textbf{Output} 页展开 \texttt{runs/seg\_yolo11n-seg\_1024/}，
        点开 \texttt{val\_batch\{0,1,2\}\_pred.jpg} 与 \texttt{results.png}
        （点击文件名即可；浏览器地址栏里那条带签名的 \texttt{storage.googleapis.com}
        直链可以直接用 \texttt{curl}/脚本下载）。
  \item 每张 \texttt{val\_batch*\_pred.jpg} 是 $3\times3$ 九宫格，每格一张验证帧；
        脚本按 $(r,c)$ 坐标裁出整格，不做任何重绘。
  \item 四格拼成 $2\times2$、加 (a)--(d) 角标，导出 PNG（$1280\times1280$）与 PDF。
\end{enumerate}
想换例子只需改脚本里的 \texttt{PANELS} 列表。若要换成"GT 对预测"的并排图，
再去 Output 页取同目录的 \texttt{val\_batch*\_labels.jpg}（GT 侧的可视化）即可。

\end{document}
